\documentclass[11pt,a4paper]{article}

\usepackage{estilo-libro}
\usetikzlibrary{matrix}
\usepackage{needspace}
% Que un título de paso no quede solo al final de una página
\let\pasobase\paso
\renewcommand{\paso}[1]{\needspace{7\baselineskip}\pasobase{#1}}

% ---------- Tablas del método húngaro (solo en este documento) ----------
% Las líneas que tapan ceros se dibujan en una capa de fondo para que los números se sigan leyendo.
\pgfdeclarelayer{fondo}
\pgfsetlayers{fondo,main}
\tikzset{
  hung/.style={matrix of nodes, row sep=0pt, column sep=0pt,
    nodes={minimum height=5.8mm, minimum width=#1, anchor=center, inner sep=1pt, font=\small},
    row 1/.style={nodes={font=\footnotesize}},
    column 1/.style={anchor=east, nodes={minimum width=0pt, inner sep=3pt}}},
  hung/.default=9mm,
  chico/.style={column 1/.append style={nodes={font=\scriptsize}}},
  tapa/.style={line width=5pt, gray!45, line cap=butt},
  asig/.style={draw, thick, circle, minimum size=5.6mm, inner sep=0pt},
  ag/.style={draw, thick, rounded corners=2pt, minimum width=21mm, minimum height=6.5mm, font=\small, fill=white},
  libre/.style={gray!55},
  elegido/.style={line width=1.8pt},
}
% \marco{matriz}{última fila}{última columna}: cuadrícula de los datos (la fila 1 y la columna 1 son nombres)
\newcommand{\marco}[3]{%
  \foreach \i in {2,...,#2} \foreach \j in {2,...,#3}
    \draw[gray!60, very thin] (#1-\i-\j.north west) rectangle (#1-\i-\j.south east);
  \draw[semithick] (#1-2-2.north west) rectangle (#1-#2-#3.south east);}
% \tfila{matriz}{fila}{última columna} y \tcol{matriz}{columna}{última fila}: líneas que tapan ceros
\newcommand{\tfila}[3]{\draw[tapa] (#1-#2-2.west) -- (#1-#2-#3.east);}
\newcommand{\tcol}[3]{\draw[tapa] (#1-2-#2.north) -- (#1-#3-#2.south);}
% \pie{matriz}{texto}: leyenda debajo de una tabla
\newcommand{\pie}[2]{\node[below=1.5mm of #1, font=\scriptsize\itshape] {#2};}
% \cero{celda}: cero asignado, encerrado en un círculo
\newcommand{\cero}[1]{\node[asig] at (#1) {};}
\newcommand{\flecha}[3]{\draw[-{Stealth[length=2.4mm]}, thick, shorten >=1mm] (#1.east) -- node[above, font=\scriptsize, align=center] {#3} (#1.east -| #2.west);}

\title{\textbf{El problema de asignación}\\[4pt]\large El método de Kuhn-Munkres (método húngaro) y ejercicios resueltos}
\author{Investigación de Operaciones · Asignación y flujo en redes}
\date{}

\begin{document}
\maketitle
\tableofcontents
\newpage

% =====================================================================
\part{¿En qué consiste el problema de asignación?}

\section{¿Qué tipo de problema resuelve?}

Muchas veces hay que \textbf{repartir trabajos entre personas} (o máquinas, o vehículos) de modo que
\textbf{cada uno haga exactamente un trabajo} y \textbf{cada trabajo lo haga exactamente uno}. Como no todos
son igual de buenos para todo, cada pareja ``agente--tarea'' tiene su propio costo o tiempo. El
\textbf{problema de asignación} consiste en elegir las parejas para que el \textbf{total sea lo más pequeño posible}
(o lo más grande, si se trata de ganancias).

Para entenderlo usaremos un ejemplo pequeño, que en la Parte II se resuelve completo:

\begin{ojo}
Un papá tiene tres hijos, \textbf{John}, \textbf{Karen} y \textbf{Terri}, y tres tareas de la casa:
\textbf{podar} el césped, \textbf{pintar} la puerta del garaje y \textbf{lavar} los autos. Cada hijo dice cuánto
cobraría por cada tarea (en dólares):
\begin{center}\small
\begin{tabular}{lccc}
 & Podar & Pintar & Lavar \\
\midrule
John  & 15 & 10 & 9 \\
Karen & 9  & 15 & 10 \\
Terri & 10 & 12 & 8 \\
\end{tabular}
\end{center}
Cada hijo hace una sola tarea. ¿A quién le toca cada una para que el papá pague lo menos posible?
\end{ojo}

Todo problema de asignación tiene estas partes:

\begin{center}\small
\begin{tabular}{p{3.2cm}p{6.4cm}p{4.6cm}}
\toprule
\textbf{Parte} & \textbf{Qué es} & \textbf{En el ejemplo} \\
\midrule
Agentes & Quienes hacen el trabajo (personas, máquinas, camiones\dots). & John, Karen y Terri \\
Tareas & Lo que hay que hacer. & Podar, pintar y lavar \\
Tabla de costos & Lo que cuesta (o tarda) cada agente en cada tarea. & La tabla de arriba \\
Regla & Cada agente hace una tarea y cada tarea la hace un agente. & Un hijo por tarea \\
Objetivo & Hacer el total lo más pequeño (o lo más grande) posible. & Pagar lo menos posible \\
\bottomrule
\end{tabular}
\end{center}

\subsection*{¿Por qué no probar todas las combinaciones?}
Con tres hijos sí se puede: John tiene 3 opciones, Karen las 2 que quedan y Terri la última, así que hay
$3 \cdot 2 \cdot 1 = 3! = 6$ formas de repartir. Estas son todas:

\begin{center}\small
\begin{tabular}{lccc}
\toprule
\textbf{John -- Karen -- Terri} & \textbf{Costo} & \textbf{John -- Karen -- Terri} & \textbf{Costo} \\
\midrule
Podar -- Pintar -- Lavar & $15 + 15 + 8 = 38$ & Pintar -- Lavar -- Podar & $10 + 10 + 10 = 30$ \\
Podar -- Lavar -- Pintar & $15 + 10 + 12 = 37$ & Lavar -- Podar -- Pintar & $9 + 9 + 12 = 30$ \\
Pintar -- Podar -- Lavar & $10 + 9 + 8 = \mathbf{27}$ & Lavar -- Pintar -- Podar & $9 + 15 + 10 = 34$ \\
\bottomrule
\end{tabular}
\end{center}

El problema es que esa cantidad crece rapidísimo. Con $n$ agentes hay $n! = n\cdot(n-1)\cdots 2\cdot 1$
formas de repartir:
\[
4! = 24, \qquad 5! = 120, \qquad 10! = 3\,628\,800, \qquad 20! \approx 2.4 \times 10^{18}.
\]
Con 20 trabajadores, una computadora que revisara mil millones de repartos por segundo tardaría unos
\textbf{77 años}. Hace falta un método que llegue a la mejor respuesta sin revisarlas todas.

% =====================================================================
\section{La idea del método húngaro}

En 1955 el matemático estadounidense \textbf{Harold Kuhn} publicó un método para resolver este problema a mano.
Lo llamó \textbf{método húngaro} porque se basa en ideas de dos matemáticos húngaros, \textbf{Dénes König} y
\textbf{Jenő Egerváry}, de los años treinta. En 1957 \textbf{James Munkres} lo revisó y demostró que siempre
termina en un tiempo razonable; por eso también se le llama \textbf{algoritmo de Kuhn-Munkres}.

La idea se apoya en una observación muy sencilla. Pase lo que pase, a John hay que pagarle \textbf{por lo menos
9 dólares}, porque su tarea más barata cuesta 9. Si a toda su fila le quitamos esos 9, lo que queda es
\textbf{cuánto más cara} sale cada tarea comparada con su mejor opción:

\begin{center}
\begin{tikzpicture}
  \matrix (a) [hung=12mm] {
          & Podar & Pintar & Lavar \\
    John  & 15 & 10 & 9 \\
  };
  \matrix (b) [hung=12mm, right=20mm of a] {
          & Podar & Pintar & Lavar \\
    John  & 6 & 1 & 0 \\
  };
  \marco{a}{2}{4} \marco{b}{2}{4}
  \flecha{a}{b}{restar 9\\a toda la fila}
  \cero{b-2-4}
\end{tikzpicture}
\end{center}

Restar la misma cantidad a \textbf{toda una fila} (o a toda una columna) \textbf{no cambia cuál es el mejor reparto}:
cualquier reparto usa exactamente una casilla de esa fila, así que \textbf{todos} los repartos bajan los mismos
9 dólares. El orden entre ellos sigue igual. Lo único que cambia es el total, que baja en esa cantidad.

El método húngaro usa esto para \textbf{fabricar ceros}. Va restando hasta que la tabla tenga muchos ceros y,
cuando se puede elegir \textbf{un cero en cada fila y en cada columna}, ese reparto es el mejor: en la tabla nueva
cuesta 0 y ningún reparto puede costar menos que 0, porque ya no quedan números negativos.

\noindent\begin{minipage}[c]{0.52\textwidth}\centering
\begin{tikzpicture}[>=Stealth]
  \foreach \n/\y in {John/0, Karen/-1.1, Terri/-2.2} \node[ag] (\n) at (0,\y) {\n};
  \foreach \n/\y in {Podar/0, Pintar/-1.1, Lavar/-2.2} \node[ag] (\n) at (5,\y) {\n};
  \foreach \a in {John, Karen, Terri} \foreach \t in {Podar, Pintar, Lavar}
    \draw[libre] (\a.east) -- (\t.west);
  \draw[elegido] (John.east) -- (Pintar.west); \node[right=1mm of Pintar, font=\scriptsize] {10};
  \draw[elegido] (Karen.east) -- (Podar.west); \node[right=1mm of Podar, font=\scriptsize] {9};
  \draw[elegido] (Terri.east) -- (Lavar.west); \node[right=1mm of Lavar, font=\scriptsize] {8};
  \node[font=\small\bfseries, above=2mm of John] {Agentes};
  \node[font=\small\bfseries, above=2mm of Podar] {Tareas};
\end{tikzpicture}
\end{minipage}\hfill
\begin{minipage}[c]{0.44\textwidth}
El dibujo muestra el problema como \textbf{parejas}: cada línea gris es una pareja posible y las gruesas son las
que elige el método en el ejemplo. Como cada agente y cada tarea tocan una sola línea gruesa, se cumple la regla.
\end{minipage}
\par

% =====================================================================
\section{Antes de empezar: preparar el problema}

El método húngaro, tal como se explica aquí, necesita tres cosas:

\begin{description}[leftmargin=0pt, itemsep=5pt]
  \item[Una tabla cuadrada.] Debe haber \textbf{tantos agentes como tareas} (se dice que el problema está
        \textbf{balanceado}). Si sobran agentes o tareas, se agregan filas o columnas ``ficticias'' con costo 0.
  \item[Minimizar.] El método busca el total \textbf{más pequeño}. Si los números son ganancias y se quiere el
        total más grande, primero se convierten en ``lo que se deja de ganar''.
  \item[Todas las parejas permitidas.] Si una pareja no se puede (un trabajador que no sabe usar una máquina),
        se le pone un costo enorme, $M$, para que el método nunca la escoja.
\end{description}

Estos tres arreglos son los \textbf{casos particulares} y se explican con detalle, con sus propios ejercicios, en el
documento \emph{Casos particulares del problema de asignación}. Los tres ejercicios de este documento ya vienen cuadrados, son de minimizar y no tienen
parejas prohibidas, así que \textbf{no hace falta preparar nada}.

% =====================================================================
\section{Las palabras que se usan}

\begin{center}\small
\begin{tabular}{p{3.6cm}p{10.6cm}}
\toprule
\textbf{Término} & \textbf{Qué significa} \\
\midrule
Agente & Quien hace un trabajo. En la tabla, cada \textbf{fila} es un agente. \\
Tarea & El trabajo que se reparte. En la tabla, cada \textbf{columna} es una tarea. \\
Tabla (o matriz) de costos & Los números $c_{ij}$: lo que cuesta que el agente $i$ haga la tarea $j$. \\
Asignación & Un reparto completo: una tarea para cada agente, sin repetir tareas. \\
Balanceado & Hay el mismo número de agentes que de tareas (la tabla es cuadrada). \\
Menor de fila / de columna & El número más pequeño de esa fila o columna. Es lo que se resta en los pasos 1 y 2. \\
Línea & Una raya que tapa una fila o una columna completa de la tabla. \\
Cero tapado / sin tapar & Un cero que queda debajo de alguna línea, o que no tiene ninguna encima. \\
Cruce & Una casilla tapada por \textbf{dos} líneas a la vez (una de fila y una de columna). \\
$k$ & El número más pequeño que quedó sin tapar. Se usa en el paso 4 para crear ceros nuevos. \\
Cero asignado & El cero que se escoge para formar una pareja. Se dibuja dentro de un círculo: \asg{0}. \\
$Z^*$ & El costo total de la mejor asignación, calculado con los números \textbf{originales}. \\
\bottomrule
\end{tabular}
\end{center}

% =====================================================================
\section{Los pasos del método húngaro}\label{sec:algoritmo}

\noindent\begin{minipage}[c]{0.64\textwidth}
Para que los pasos se entiendan, debajo de cada uno se muestra lo que ocurre en el ejercicio 2 de la Parte II
(\textbf{Fowle}: tres líderes de proyecto, Terry, Carle y McClymonds, y tres clientes, con los días que tardaría cada
uno). Se usa ese ejemplo porque necesita los cinco pasos.
\end{minipage}\hfill
\begin{minipage}[c]{0.32\textwidth}\centering\small
$\begin{array}{l|ccc}
 & \text{C1} & \text{C2} & \text{C3}\\ \hline
\text{Terry} & 10 & 15 & 9\\
\text{Carle} & 9 & 18 & 5\\
\text{McClymonds} & 6 & 14 & 3
\end{array}$
\end{minipage}
\par

\paso{Paso 1. Restar el menor de cada fila}
A cada número de una fila se le resta el número más pequeño de esa fila. Así cada fila queda con al menos un cero.

\textit{En el ejemplo:} los menores son 9, 5 y 3. La fila de Terry $(10,\ 15,\ 9)$ queda $(1,\ 6,\ 0)$; la de Carle,
$(4,\ 13,\ 0)$; y la de McClymonds, $(3,\ 11,\ 0)$.

\paso{Paso 2. Restar el menor de cada columna}
Ahora se hace lo mismo por columnas, sobre la tabla que dejó el paso 1. Así cada columna queda con al menos un cero.

\textit{En el ejemplo:} los menores de las columnas son 1, 6 y 0. La tabla queda con filas $(0,\ 0,\ 0)$,
$(3,\ 7,\ 0)$ y $(2,\ 5,\ 0)$.

\paso{Paso 3. Tapar todos los ceros con el menor número de líneas}
Se trazan líneas sobre filas o columnas completas hasta que no quede ningún cero sin tapar, usando
\textbf{las menos líneas posibles}. Una forma práctica es tapar primero la fila o columna que tenga más ceros
(esta regla no siempre da el mínimo: si después no se logra asignar, hay que revisar las líneas).
\begin{itemize}[itemsep=2pt]
  \item Si se necesitan \textbf{tantas líneas como filas} ($n$), ya se puede asignar: se salta al paso 5.
  \item Si alcanzan \textbf{menos de $n$ líneas}, todavía no hay un cero para cada fila: se pasa al paso 4.
\end{itemize}
El número mínimo de líneas es justo la cantidad de parejas que se pueden formar con ceros sin repetir fila ni
columna (lo demostró König). Por eso, con $n$ líneas, siempre hay una asignación completa.

\textit{En el ejemplo:} basta con 2 líneas, la fila de Terry y la columna C3. Como $2 < 3$, todavía no se puede
asignar.

\paso{Paso 4. Ajustar con el menor número sin tapar}
Se busca $k$, el número más pequeño que \textbf{no} tiene ninguna línea encima. Entonces:
\begin{enumerate}[itemsep=2pt]
  \item a todos los números \textbf{sin tapar} se les resta $k$ (así aparece al menos un cero nuevo);
  \item a los \textbf{cruces} de dos líneas se les suma $k$;
  \item los números tapados por una sola línea \textbf{no cambian}.
\end{enumerate}
Esto equivale a restar $k$ a las filas sin línea y sumar $k$ a las columnas con línea, así que sigue siendo
``restar a filas y columnas completas'' y no cambia cuál es el mejor reparto. Después se vuelve al paso 3.

\textit{En el ejemplo:}

\begin{center}
\begin{tikzpicture}
  \matrix (a) [hung] {
               & C1 & C2 & C3 \\
    Terry      & 0 & 0 & 0 \\
    Carle      & 3 & 7 & 0 \\
    McClymonds & 2 & 5 & 0 \\
  };
  \matrix (b) [hung, right=22mm of a] {
               & C1 & C2 & C3 \\
    Terry      & 0 & 0 & 2 \\
    Carle      & 1 & 5 & 0 \\
    McClymonds & 0 & 3 & 0 \\
  };
  \begin{pgfonlayer}{fondo}
    \marco{a}{4}{4} \marco{b}{4}{4}
    \tfila{a}{2}{4} \tcol{a}{4}{4}
  \end{pgfonlayer}
  \flecha{a}{b}{$k = 2$}
\end{tikzpicture}
\end{center}
Sin tapar quedaron 3, 7, 2 y 5; el menor es $k = 2$. Esos cuatro bajan 2 (aparece el cero de McClymonds en C1),
el cruce de la fila de Terry con la columna C3 sube de 0 a 2, y el resto queda igual. Al volver al paso 3 ya hacen
falta 3 líneas.

\paso{Paso 5. Asignar un cero por fila y por columna}
Se empieza por una fila (o columna) que tenga \textbf{un solo cero}: ese cero se encierra en un círculo y se tachan
su fila y su columna, porque esa tarea y ese agente ya están ocupados. Se repite con lo que queda. El costo total se
calcula sumando los números \textbf{originales} de las casillas elegidas.

\textit{En el ejemplo:} Carle solo tiene un cero (C3); después McClymonds solo tiene libre el de C1, y a Terry le
queda C2. El total es $15 + 5 + 6 = 26$ días.

\needspace{13\baselineskip}
\begin{idea}[Resumen de los pasos]
\begin{enumerate}[leftmargin=*, itemsep=3pt]
  \item[\textbf{0.}] Preparar: tabla cuadrada, de minimizar y sin parejas prohibidas.
  \item \textbf{Filas:} a cada fila se le resta su menor.
  \item \textbf{Columnas:} a cada columna se le resta su menor.
  \item \textbf{Líneas:} se tapan todos los ceros con el menor número de líneas. Si son $n$, se va al paso 5.
  \item \textbf{Ajuste:} $k$ = menor número sin tapar; se resta $k$ a lo no tapado y se suma $k$ a los cruces.
        Se vuelve al paso 3.
  \item \textbf{Asignar:} un cero por fila y por columna, empezando por las filas o columnas con un solo cero.
        El total se calcula con los costos originales.
\end{enumerate}
\end{idea}

\begin{center}
\begin{tikzpicture}[node distance=6mm and 9mm, >=Stealth,
  caja/.style={draw, rounded corners=2pt, fill=gray!8, align=center, minimum width=5.6cm, inner sep=5pt, font=\small},
  decision/.style={draw, diamond, aspect=2.4, fill=gray!18, align=center, inner sep=1pt, font=\small},
  fin/.style={draw, rounded corners=8pt, fill=gray!25, align=center, minimum width=2.8cm, inner sep=6pt, font=\small}]
  \node[caja] (ini) {Preparar la tabla\\ (cuadrada y de minimizar)};
  \node[caja, below=of ini] (fil) {1. Restar el menor de cada fila};
  \node[caja, below=of fil] (col) {2. Restar el menor de cada columna};
  \node[caja, below=of col] (lin) {3. Tapar los ceros con\\ el menor número de líneas};
  \node[decision, below=of lin] (dec) {¿Hay $n$ líneas?};
  \node[caja, below=of dec] (asg) {5. Asignar un cero\\ por fila y por columna};
  \node[fin, below=of asg] (opt) {\textbf{Asignación óptima}};
  \node[caja, minimum width=3.6cm, right=of dec] (aju) {4. Ajustar con $k$:\\ $-k$ sin tapar, $+k$ en cruces};
  \draw[->] (ini) -- (fil);
  \draw[->] (fil) -- (col);
  \draw[->] (col) -- (lin);
  \draw[->] (lin) -- (dec);
  \draw[->] (dec) -- node[right, font=\scriptsize] {sí} (asg);
  \draw[->] (asg) -- (opt);
  \draw[->] (dec) -- node[above, font=\scriptsize] {no} (aju);
  \draw[->] (aju.north) |- node[pos=0.25, right, font=\scriptsize] {otra vez} (lin.east);
\end{tikzpicture}
\end{center}

% =====================================================================
\section{¿Cómo se sabe que terminó?}

\begin{description}[leftmargin=0pt, itemsep=5pt]
  \item[Solución óptima.] En el paso 3 hacen falta $n$ líneas para tapar los ceros. Entonces se puede elegir un cero
        en cada fila y en cada columna, y ese reparto es el mejor.
  \item[Siempre termina.] Cada ajuste del paso 4 hace bajar la suma de la tabla, y la tabla nunca tiene números
        negativos, así que no puede bajar para siempre. En la práctica bastan muy pocos ajustes.
  \item[Varias soluciones óptimas.] Si al asignar se puede elegir más de un juego de ceros (por ejemplo, dos filas con
        dos ceros en las mismas columnas), hay más de un reparto con el mismo costo mínimo.
  \item[Comprobación rápida.] La suma de todo lo que se restó (los menores de filas y columnas, y en cada ajuste
        $k \times$ (filas sin línea $-$ columnas con línea)) es igual al costo mínimo. En el ejemplo: filas $9+5+3 = 17$, columnas $1+6+0 = 7$ y el ajuste
        $2 \times (2 - 1) = 2$; en total $17 + 7 + 2 = 26$.
\end{description}

% =====================================================================
\section{Relación con el transporte y la programación lineal}

El problema de asignación es un \textbf{caso especial del problema de transporte}: cada agente es un ``origen'' que
ofrece 1 unidad de trabajo y cada tarea es un ``destino'' que pide 1 unidad. También se puede escribir como un
problema de programación lineal. Se usa una variable por pareja:
\[
x_{ij} = \begin{cases} 1 & \text{si el agente } i \text{ hace la tarea } j,\\ 0 & \text{si no.}\end{cases}
\]
\[
\begin{aligned}
\text{Min } Z &= \textstyle\sum_{i}\sum_{j} c_{ij}\, x_{ij}\\
\text{s.a. } \quad \textstyle\sum_{j} x_{ij} &= 1 \quad \text{para cada agente } i \nota{cada agente hace una tarea}\\
\textstyle\sum_{i} x_{ij} &= 1 \quad \text{para cada tarea } j \nota{cada tarea la hace un agente}\\
x_{ij} &\in \{0, 1\}
\end{aligned}
\]
Una propiedad muy útil: si se resuelve con el simplex permitiendo cualquier $x_{ij} \ge 0$, la respuesta sale
\textbf{sola con ceros y unos}. Aun así, con $n$ agentes el simplex maneja $n^2$ variables y $2n$ restricciones,
mientras que el método húngaro trabaja directo sobre la tabla de $n \times n$, que es mucho más cómodo.

% =====================================================================
\section{Ventajas y desventajas}

\begin{center}\small
\begin{tabular}{p{7.2cm}p{7.2cm}}
\toprule
\textbf{Ventajas} & \textbf{Desventajas} \\
\midrule
\begin{itemize}[leftmargin=*, itemsep=2pt, topsep=2pt]
  \item Siempre encuentra la mejor asignación, sin probar las $n!$ combinaciones.
  \item Solo usa restas y sumas: se hace a mano sobre la misma tabla.
  \item Es rápido también en computadora (tarda del orden de $n^3$ operaciones).
  \item Lo restado da una comprobación del costo mínimo.
\end{itemize}
&
\begin{itemize}[leftmargin=*, itemsep=2pt, topsep=2pt]
  \item Encontrar a ojo el menor número de líneas puede ser difícil en tablas grandes.
  \item Necesita una tabla cuadrada y de minimizar; los demás casos hay que prepararlos antes.
  \item Solo sirve si cada agente hace \textbf{una} tarea (si hace varias, hay que duplicar filas).
\end{itemize}
\\
\bottomrule
\end{tabular}
\end{center}

% =====================================================================
\part{Ejercicios resueltos}
\setcounter{section}{0}

Los tres ejercicios siguientes aplican, uno por uno, los pasos de la sección~\ref{sec:algoritmo}. En las tablas, las
\textbf{franjas grises} son las líneas que tapan ceros y los \textbf{círculos} marcan los ceros asignados. Al final de
cada ejercicio se dibujan las parejas elegidas.

% =====================================================================
\section{Los hijos de Joe Klyne: tareas de la casa}
\textit{Fuente: Taha, Investigación de Operaciones, capítulo del modelo de transporte y sus variantes (modelo de asignación).}

\subsection*{Planteamiento del problema}
Joe Klyne quiere que sus tres hijos, John, Karen y Terri, ganen algo de dinero. Escoge tres tareas: podar el
césped, pintar la puerta del garaje y lavar los autos de la familia. Para que no haya peleas, les pide que digan
cuánto cobrarían por cada tarea:

\begin{center}\small
\begin{tabular}{lccc}
\toprule
 & \textbf{Podar} & \textbf{Pintar} & \textbf{Lavar} \\
\midrule
John  & \$15 & \$10 & \$9 \\
Karen & \$9  & \$15 & \$10 \\
Terri & \$10 & \$12 & \$8 \\
\bottomrule
\end{tabular}
\end{center}

La pregunta: \textbf{¿qué tarea le toca a cada hijo para que Joe pague lo menos posible?}

\subsection*{Formulación del modelo}
\begin{itemize}[itemsep=2pt]
  \item $x_{ij} = 1$ si el hijo $i$ hace la tarea $j$, y $0$ si no (filas: 1 John, 2 Karen, 3 Terri; columnas:
        1 podar, 2 pintar, 3 lavar).
  \item \textbf{Lo que paga Joe:} $Z = 15x_{11} + 10x_{12} + 9x_{13} + 9x_{21} + 15x_{22} + 10x_{23} + 10x_{31} + 12x_{32} + 8x_{33}$.
  \item \textbf{Cada hijo hace una tarea:} $x_{11} + x_{12} + x_{13} = 1$, y lo mismo para Karen y Terri.
  \item \textbf{Cada tarea la hace un hijo:} $x_{11} + x_{21} + x_{31} = 1$, y lo mismo para pintar y lavar.
\end{itemize}
\textbf{Preparación:} hay 3 hijos y 3 tareas (la tabla es cuadrada) y se quiere pagar lo mínimo, así que se puede
empezar directamente.

\paso{Paso 1. Restar el menor de cada fila}
\begin{center}
\begin{tikzpicture}
  \matrix (a) [hung=10mm] {
          & Podar & Pintar & Lavar & \footnotesize menor \\
    John & 15 & 10 & 9 & \footnotesize 9 \\
    Karen & 9 & 15 & 10 & \footnotesize 9 \\
    Terri & 10 & 12 & 8 & \footnotesize 8 \\
  };
  \matrix (b) [hung=10mm, anchor=north west] at ([xshift=10mm]a.north east) {
          & Podar & Pintar & Lavar \\
    John & 6 & 1 & 0 \\
    Karen & 0 & 6 & 1 \\
    Terri & 2 & 4 & 0 \\
    \footnotesize menor & \footnotesize 0 & \footnotesize 1 & \footnotesize 0 \\
  };
  \matrix (c) [hung=10mm, anchor=north west] at ([xshift=10mm]b.north east) {
          & Podar & Pintar & Lavar \\
    John & 6 & 0 & 0 \\
    Karen & 0 & 5 & 1 \\
    Terri & 2 & 3 & 0 \\
  };
  \marco{a}{4}{4} \marco{b}{4}{4} \marco{c}{4}{4}
  \flecha{a}{b}{paso 1} \flecha{b}{c}{paso 2}
  \pie{a}{tabla original} \pie{b}{después del paso 1} \pie{c}{después del paso 2}
\end{tikzpicture}
\end{center}
Los menores son 9, 9 y 8 (columna ``menor'' de la tabla original). Por ejemplo,
$\text{John:}\ (15,\ 10,\ 9) - 9 = (6,\ 1,\ 0)$: pase lo que pase, John cobra al menos 9.
Igual con Karen ($-9$) y Terri ($-8$).

\paso{Paso 2. Restar el menor de cada columna}
Podar y lavar ya tenían un cero (su menor es 0, no cambian). En \textbf{pintar} el menor es 1:
\[
(1,\ 6,\ 4) - 1 = (0,\ 5,\ 3) \nota{aparece un cero para John en pintar}
\]

\paso{Paso 3. Tapar los ceros con el menor número de líneas}
\begin{center}
\begin{tikzpicture}
  \matrix (a) [hung=12mm] {
          & Podar & Pintar & Lavar \\
    John  & 6 & 0 & 0 \\
    Karen & 0 & 5 & 1 \\
    Terri & 2 & 3 & 0 \\
  };
  \matrix (b) [hung=12mm, right=22mm of a] {
          & Podar & Pintar & Lavar \\
    John  & 6 & 0 & 0 \\
    Karen & 0 & 5 & 1 \\
    Terri & 2 & 3 & 0 \\
  };
  \begin{pgfonlayer}{fondo}
    \marco{a}{4}{4} \marco{b}{4}{4}
    \tfila{a}{2}{4} \tfila{a}{3}{4} \tfila{a}{4}{4}
  \end{pgfonlayer}
  \cero{b-3-2} \cero{b-4-4} \cero{b-2-3}
  \flecha{a}{b}{asignar}
  \pie{a}{Paso 3: 3 líneas} \pie{b}{Paso 5: ceros asignados}
\end{tikzpicture}
\end{center}
Los ceros están en (John, pintar), (John, lavar), (Karen, podar) y (Terri, lavar). No alcanza con 2 líneas: el cero
de Karen en podar necesita una línea propia, y los otros tres necesitan dos más, porque (John, pintar) y
(Terri, lavar) no comparten ni fila ni columna. Hacen falta \textbf{3 líneas} (aquí, las tres filas).
\[
3 \text{ líneas} = 3 \text{ filas} \nota{ya se puede asignar: no hace falta el paso 4}
\]

\paso{Paso 5. Asignar un cero por fila y por columna}
En la tabla de la derecha:
\begin{align*}
1)\quad & \text{Karen} \to \text{Podar} \nota{único cero de su fila; se tachan su fila y la columna podar}\\
2)\quad & \text{Terri} \to \text{Lavar} \nota{único cero de su fila; se tachan su fila y la columna lavar}\\
3)\quad & \text{John} \to \text{Pintar} \nota{su otro cero (lavar) ya está tachado}
\end{align*}

\subsection*{Resultado}
\[
\res{\text{John} \to \text{Pintar},\ \ \text{Karen} \to \text{Podar},\ \ \text{Terri} \to \text{Lavar},\quad Z^* = 27}
\]
\begin{itemize}[itemsep=2pt]
  \item John pinta la puerta (\$10), Karen poda el césped (\$9) y Terri lava los autos (\$8).
  \item Joe paga en total $10 + 9 + 8 = \mathbf{\$27}$, con los precios originales.
  \item Karen y Terri hacen su tarea más barata. John cobraría menos por lavar (\$9), pero esa tarea le sale más
        barata a Terri (\$8), y John es quien menos cobra por pintar (\$10).
\end{itemize}

\textbf{Comprobación:} cada hijo tiene una tarea y cada tarea un hijo \checkmark. Lo restado suma
$(9 + 9 + 8) + (0 + 1 + 0) = 27$, igual al costo \checkmark. Y en la tabla de las 6 combinaciones de la Parte I,
el menor costo es justo 27, el de este reparto \checkmark.

\subsection*{Interpretación gráfica}
\noindent\begin{minipage}[c]{0.52\textwidth}\centering
\begin{tikzpicture}[>=Stealth]
  \foreach \n/\y in {John/0, Karen/-0.95, Terri/-1.9} \node[ag] (\n) at (0,\y) {\n};
  \foreach \n/\y in {Podar/0, Pintar/-0.95, Lavar/-1.9} \node[ag] (\n) at (5,\y) {\n};
  \foreach \a in {John, Karen, Terri} \foreach \t in {Podar, Pintar, Lavar}
    \draw[libre] (\a.east) -- (\t.west);
  \draw[elegido] (John.east) -- (Pintar.west); \node[right=1mm of Pintar, font=\scriptsize] {\$10};
  \draw[elegido] (Karen.east) -- (Podar.west); \node[right=1mm of Podar, font=\scriptsize] {\$9};
  \draw[elegido] (Terri.east) -- (Lavar.west); \node[right=1mm of Lavar, font=\scriptsize] {\$8};
  \node[font=\small\bfseries, above=2mm of John] {Hijos};
  \node[font=\small\bfseries, above=2mm of Podar] {Tareas};
\end{tikzpicture}
\end{minipage}\hfill
\begin{minipage}[c]{0.44\textwidth}
De las 9 parejas posibles (líneas grises) se usan 3 (gruesas), una por hijo y una por tarea. En palabras: a Joe le
conviene que \textbf{John pinte, Karen pode y Terri lave}; así paga \$27, que es lo menos posible.
\end{minipage}
\par

% =====================================================================
\section{Fowle Marketing Research: líderes de proyecto y clientes}
\textit{Fuente: Anderson, Sweeney y Williams, Métodos cuantitativos para los negocios, capítulo del problema de asignación.}

\subsection*{Planteamiento del problema}
Fowle Marketing Research acaba de recibir pedidos de estudios de mercado de \textbf{tres clientes nuevos}. Tiene
\textbf{tres líderes de proyecto} libres: Terry, Carle y McClymonds. Cada líder trabajaría con un solo cliente. Según
su experiencia, los días que tardaría cada uno en terminar cada estudio son:

\begin{center}\small
\begin{tabular}{lccc}
\toprule
\textbf{Líder de proyecto} & \textbf{Cliente 1 (C1)} & \textbf{Cliente 2 (C2)} & \textbf{Cliente 3 (C3)} \\
\midrule
Terry      & 10 & 15 & 9 \\
Carle      & 9  & 18 & 5 \\
McClymonds & 6  & 14 & 3 \\
\bottomrule
\end{tabular}
\end{center}

La pregunta: \textbf{¿qué líder conviene asignar a cada cliente para que el total de días sea el menor?}

\subsection*{Formulación del modelo}
\begin{itemize}[itemsep=2pt]
  \item $x_{ij} = 1$ si el líder $i$ (1 Terry, 2 Carle, 3 McClymonds) atiende al cliente $j$, y $0$ si no.
  \item \textbf{Días en total:} $Z = 10x_{11} + 15x_{12} + 9x_{13} + 9x_{21} + 18x_{22} + 5x_{23} + 6x_{31} + 14x_{32} + 3x_{33}$.
  \item \textbf{Cada líder atiende un cliente:} $x_{i1} + x_{i2} + x_{i3} = 1$ para $i = 1, 2, 3$.
  \item \textbf{Cada cliente tiene un líder:} $x_{1j} + x_{2j} + x_{3j} = 1$ para $j = 1, 2, 3$.
\end{itemize}
\textbf{Preparación:} 3 líderes y 3 clientes, de minimizar: no hace falta preparar nada.

\paso{Paso 1. Restar el menor de cada fila}
\begin{center}
\begin{tikzpicture}
  \matrix (a) [hung=8mm, chico] {
          & C1 & C2 & C3 & \footnotesize menor \\
    Terry & 10 & 15 & 9 & \footnotesize 9 \\
    Carle & 9 & 18 & 5 & \footnotesize 5 \\
    McClymonds & 6 & 14 & 3 & \footnotesize 3 \\
  };
  \matrix (b) [hung=8mm, chico, anchor=north west] at ([xshift=8.5mm]a.north east) {
          & C1 & C2 & C3 \\
    Terry & 1 & 6 & 0 \\
    Carle & 4 & 13 & 0 \\
    McClymonds & 3 & 11 & 0 \\
    \footnotesize menor & \footnotesize 1 & \footnotesize 6 & \footnotesize 0 \\
  };
  \matrix (c) [hung=8mm, chico, anchor=north west] at ([xshift=8.5mm]b.north east) {
          & C1 & C2 & C3 \\
    Terry & 0 & 0 & 0 \\
    Carle & 3 & 7 & 0 \\
    McClymonds & 2 & 5 & 0 \\
  };
  \marco{a}{4}{4} \marco{b}{4}{4} \marco{c}{4}{4}
  \flecha{a}{b}{paso 1} \flecha{b}{c}{paso 2}
  \pie{a}{tabla original} \pie{b}{después del paso 1} \pie{c}{después del paso 2}
\end{tikzpicture}
\end{center}
Los menores son 9, 5 y 3, y los tres están en la columna C3: \textbf{los tres líderes son más rápidos con el
cliente 3}. Por ejemplo, $\text{Terry:}\ (10,\ 15,\ 9) - 9 = (1,\ 6,\ 0)$.

\paso{Paso 2. Restar el menor de cada columna}
C1 baja 1: $(1,\ 4,\ 3) - 1 = (0,\ 3,\ 2)$. C2 baja 6: $(6,\ 13,\ 11) - 6 = (0,\ 7,\ 5)$, porque el cliente 2
le cuesta a cualquiera al menos 6 días más que su mejor opción. C3 ya tenía ceros y no cambia.

\paso{Paso 3. Tapar los ceros con el menor número de líneas}
\begin{center}
\begin{tikzpicture}
  \matrix (a) [hung] {
               & C1 & C2 & C3 \\
    Terry      & 0 & 0 & 0 \\
    Carle      & 3 & 7 & 0 \\
    McClymonds & 2 & 5 & 0 \\
  };
  \matrix (b) [hung, right=22mm of a] {
               & C1 & C2 & C3 \\
    Terry      & 0 & 0 & 2 \\
    Carle      & 1 & 5 & 0 \\
    McClymonds & 0 & 3 & 0 \\
  };
  \begin{pgfonlayer}{fondo}
    \marco{a}{4}{4} \marco{b}{4}{4}
    \tfila{a}{2}{4} \tcol{a}{4}{4}
  \end{pgfonlayer}
  \flecha{a}{b}{$k = 2$}
  \pie{a}{Paso 3: 2 líneas} \pie{b}{Paso 4: después del ajuste}
\end{tikzpicture}
\end{center}
La fila de Terry tapa sus tres ceros y la columna C3 tapa los de Carle y McClymonds: con \textbf{2 líneas} se
tapan todos.
\[
2 \text{ líneas} < 3 \text{ filas} \nota{todavía no: Carle y McClymonds solo tienen ceros en C3}
\]

\paso{Paso 4. Ajustar con el menor número sin tapar}
Sin tapar quedaron $3,\ 7,\ 2,\ 5$, así que $k = 2$.
\begin{align*}
\text{sin tapar} &: 3 - 2 = 1,\quad 7 - 2 = 5,\quad 2 - 2 = \mathbf{0},\quad 5 - 2 = 3 \nota{nace un cero}\\
\text{cruce (Terry, C3)} &: 0 + 2 = 2 \nota{tapado por dos líneas: se suma $k$}\\
\text{tapados una vez} &: \text{no cambian}
\end{align*}

\paso{Paso 3 (otra vez). Tapar los ceros}
\begin{center}
\begin{tikzpicture}
  \matrix (a) [hung] {
               & C1 & C2 & C3 \\
    Terry      & 0 & 0 & 2 \\
    Carle      & 1 & 5 & 0 \\
    McClymonds & 0 & 3 & 0 \\
  };
  \matrix (b) [hung, right=22mm of a] {
               & C1 & C2 & C3 \\
    Terry      & 0 & 0 & 2 \\
    Carle      & 1 & 5 & 0 \\
    McClymonds & 0 & 3 & 0 \\
  };
  \begin{pgfonlayer}{fondo}
    \marco{a}{4}{4} \marco{b}{4}{4}
    \tfila{a}{2}{4} \tfila{a}{3}{4} \tfila{a}{4}{4}
  \end{pgfonlayer}
  \cero{b-3-4} \cero{b-4-2} \cero{b-2-3}
  \flecha{a}{b}{asignar}
  \pie{a}{Paso 3: 3 líneas} \pie{b}{Paso 5: ceros asignados}
\end{tikzpicture}
\end{center}
Con el cero nuevo hacen falta \textbf{3 líneas} (las tres filas). Como $3 = 3$, ya se puede asignar.

\paso{Paso 5. Asignar un cero por fila y por columna}
En la tabla de la derecha:
\begin{align*}
1)\quad & \text{Carle} \to \text{C3} \nota{único cero de su fila; se tacha la columna C3}\\
2)\quad & \text{McClymonds} \to \text{C1} \nota{su cero de C3 ya está tachado: le queda C1}\\
3)\quad & \text{Terry} \to \text{C2} \nota{C1 ya está ocupado: le queda C2}
\end{align*}

\subsection*{Resultado}
\[
\res{\text{Terry} \to \text{C2},\ \ \text{Carle} \to \text{C3},\ \ \text{McClymonds} \to \text{C1},\quad Z^* = 26}
\]
\begin{itemize}[itemsep=2pt]
  \item Terry atiende al cliente 2 (15 días), Carle al cliente 3 (5 días) y McClymonds al cliente 1 (6 días).
  \item En total, con los tiempos originales: $15 + 5 + 6 = \mathbf{26}$ días.
  \item El cliente 2 es lento para todos. Comparado con su mejor cliente, Terry es quien menos días pierde con él
        (6 días más, contra 13 de Carle y 11 de McClymonds); por eso se lo queda Terry.
\end{itemize}

\textbf{Comprobación:} cada líder tiene un cliente y cada cliente un líder \checkmark. Lo restado:
filas $9 + 5 + 3 = 17$, columnas $1 + 6 + 0 = 7$; el ajuste restó $2$ a las dos filas sin línea y sumó $2$ a la
columna con línea, o sea bajó $2 \cdot 2 - 2 = 2$. Total $17 + 7 + 2 = 26$ \checkmark. Las 6 combinaciones dan
31, 29, 27, \textbf{26}, 32 y 33: el mínimo es 26 \checkmark.

\subsection*{Interpretación gráfica}
\noindent\begin{minipage}[c]{0.52\textwidth}\centering
\begin{tikzpicture}[>=Stealth]
  \foreach \n/\y in {Terry/0, Carle/-0.95, McClymonds/-1.9} \node[ag] (\n) at (0,\y) {\n};
  \foreach \n/\y in {C1/0, C2/-0.95, C3/-1.9} \node[ag] (\n) at (5,\y) {Cliente \n};
  \foreach \a in {Terry, Carle, McClymonds} \foreach \t in {C1, C2, C3}
    \draw[libre] (\a.east) -- (\t.west);
  \draw[elegido] (Terry.east) -- (C2.west); \node[right=1mm of C2, font=\scriptsize] {15 d};
  \draw[elegido] (Carle.east) -- (C3.west); \node[right=1mm of C3, font=\scriptsize] {5 d};
  \draw[elegido] (McClymonds.east) -- (C1.west); \node[right=1mm of C1, font=\scriptsize] {6 d};
  \node[font=\small\bfseries, above=2mm of Terry] {Líderes};
  \node[font=\small\bfseries, above=2mm of C1] {Clientes};
\end{tikzpicture}
\end{minipage}\hfill
\begin{minipage}[c]{0.44\textwidth}
En palabras: aunque los tres líderes son más rápidos con el cliente 3, solo uno puede quedárselo. Se lo queda
\textbf{Carle}, que es el que más ganaría con él; McClymonds toma el cliente 1 y Terry el 2. Así se ocupan
\textbf{26 días} de trabajo, lo mínimo posible.
\end{minipage}
\par

% =====================================================================
\section{Machineco: máquinas y trabajos}
\textit{Fuente: Winston, Investigación de operaciones: aplicaciones y algoritmos, capítulo de los problemas de transporte y asignación.}

\subsection*{Planteamiento del problema}
Machineco tiene \textbf{cuatro máquinas} y \textbf{cuatro trabajos} por terminar. Cada máquina debe hacer un solo
trabajo. Antes de empezar, cada máquina se tiene que preparar, y ese tiempo de preparación (en horas) depende del
trabajo:

\begin{center}\small
\begin{tabular}{lcccc}
\toprule
 & \textbf{Trabajo 1 (T1)} & \textbf{Trabajo 2 (T2)} & \textbf{Trabajo 3 (T3)} & \textbf{Trabajo 4 (T4)} \\
\midrule
Máquina 1 (M1) & 14 & 5  & 8 & 7 \\
Máquina 2 (M2) & 2  & 12 & 6 & 5 \\
Máquina 3 (M3) & 7  & 8  & 3 & 9 \\
Máquina 4 (M4) & 2  & 4  & 6 & 10 \\
\bottomrule
\end{tabular}
\end{center}

La pregunta: \textbf{¿qué trabajo debe hacer cada máquina para que el tiempo total de preparación sea el menor?}
Aquí ya hay $4! = 24$ repartos posibles.

\subsection*{Formulación del modelo}
\begin{itemize}[itemsep=2pt]
  \item $x_{ij} = 1$ si la máquina $i$ hace el trabajo $j$, y $0$ si no ($i, j = 1, \dots, 4$).
  \item \textbf{Horas de preparación:} $Z = 14x_{11} + 5x_{12} + 8x_{13} + 7x_{14} + 2x_{21} + \dots + 10x_{44}$
        (los 16 números de la tabla).
  \item \textbf{Cada máquina hace un trabajo:} $\sum_j x_{ij} = 1$; \quad \textbf{cada trabajo en una máquina:}
        $\sum_i x_{ij} = 1$.
\end{itemize}
\textbf{Preparación:} 4 máquinas y 4 trabajos, de minimizar: se empieza directamente.

\paso{Paso 1. Restar el menor de cada fila}
\begin{center}
\begin{tikzpicture}
  \matrix (a) [hung=8mm] {
          & T1 & T2 & T3 & T4 & \footnotesize menor \\
    M1 & 14 & 5 & 8 & 7 & \footnotesize 5 \\
    M2 & 2 & 12 & 6 & 5 & \footnotesize 2 \\
    M3 & 7 & 8 & 3 & 9 & \footnotesize 3 \\
    M4 & 2 & 4 & 6 & 10 & \footnotesize 2 \\
  };
  \matrix (b) [hung=8mm, anchor=north west] at ([xshift=12mm]a.north east) {
          & T1 & T2 & T3 & T4 \\
    M1 & 9 & 0 & 3 & 2 \\
    M2 & 0 & 10 & 4 & 3 \\
    M3 & 4 & 5 & 0 & 6 \\
    M4 & 0 & 2 & 4 & 8 \\
    \footnotesize menor & \footnotesize 0 & \footnotesize 0 & \footnotesize 0 & \footnotesize 2 \\
  };
  \matrix (c) [hung=8mm, anchor=north west] at ([xshift=12mm]b.north east) {
          & T1 & T2 & T3 & T4 \\
    M1 & 9 & 0 & 3 & 0 \\
    M2 & 0 & 10 & 4 & 1 \\
    M3 & 4 & 5 & 0 & 4 \\
    M4 & 0 & 2 & 4 & 6 \\
  };
  \marco{a}{5}{5} \marco{b}{5}{5} \marco{c}{5}{5}
  \flecha{a}{b}{paso 1} \flecha{b}{c}{paso 2}
  \pie{a}{tabla original} \pie{b}{después del paso 1} \pie{c}{después del paso 2}
\end{tikzpicture}
\end{center}
Los menores son 5, 2, 3 y 2. Por ejemplo, $\text{M1:}\ (14,\ 5,\ 8,\ 7) - 5 = (9,\ 0,\ 3,\ 2)$: la máquina 1
tarda al menos 5 h en prepararse, haga el trabajo que haga.

\paso{Paso 2. Restar el menor de cada columna}
T1, T2 y T3 ya tienen un cero. Solo cambia T4, cuyo menor es 2:
\[
(2,\ 3,\ 6,\ 8) - 2 = (0,\ 1,\ 4,\ 6) \nota{aparece un cero para M1 en T4}
\]

\paso{Paso 3. Tapar los ceros con el menor número de líneas}
\begin{center}
\begin{tikzpicture}
  \matrix (a) [hung] {
       & T1 & T2 & T3 & T4 \\
    M1 & 9 & 0  & 3 & 0 \\
    M2 & 0 & 10 & 4 & 1 \\
    M3 & 4 & 5  & 0 & 4 \\
    M4 & 0 & 2  & 4 & 6 \\
  };
  \matrix (b) [hung, right=22mm of a] {
       & T1 & T2 & T3 & T4 \\
    M1 & 10 & 0 & 3 & 0 \\
    M2 & 0  & 9 & 3 & 0 \\
    M3 & 5  & 5 & 0 & 4 \\
    M4 & 0  & 1 & 3 & 5 \\
  };
  \begin{pgfonlayer}{fondo}
    \marco{a}{5}{5} \marco{b}{5}{5}
    \tfila{a}{2}{5} \tfila{a}{4}{5} \tcol{a}{2}{5}
  \end{pgfonlayer}
  \flecha{a}{b}{$k = 1$}
  \pie{a}{Paso 3: 3 líneas} \pie{b}{Paso 4: después del ajuste}
\end{tikzpicture}
\end{center}
La fila M1 tapa sus dos ceros (T2 y T4), la fila M3 tapa el de T3 y la columna T1 tapa los de M2 y M4.
\[
3 \text{ líneas} < 4 \text{ filas} \nota{todavía no: M2 y M4 solo tienen ceros en T1, y no pueden compartirlo}
\]

\paso{Paso 4. Ajustar con el menor número sin tapar}
Sin tapar quedan las filas M2 y M4 sin la columna T1: $10,\ 4,\ 1$ y $2,\ 4,\ 6$. El menor es $k = 1$.
\begin{align*}
\text{M2 sin tapar} &: (10,\ 4,\ 1) - 1 = (9,\ 3,\ \mathbf{0}) \nota{nace un cero: M2 en T4}\\
\text{M4 sin tapar} &: (2,\ 4,\ 6) - 1 = (1,\ 3,\ 5)\\
\text{cruces} &: (\text{M1, T1})\ 9 + 1 = 10, \quad (\text{M3, T1})\ 4 + 1 = 5 \nota{dos líneas: se suma $k$}
\end{align*}
Los demás números (tapados por una sola línea) no cambian.

\paso{Paso 3 (otra vez). Tapar los ceros}
\begin{center}
\begin{tikzpicture}
  \matrix (a) [hung] {
       & T1 & T2 & T3 & T4 \\
    M1 & 10 & 0 & 3 & 0 \\
    M2 & 0  & 9 & 3 & 0 \\
    M3 & 5  & 5 & 0 & 4 \\
    M4 & 0  & 1 & 3 & 5 \\
  };
  \matrix (b) [hung, right=22mm of a] {
       & T1 & T2 & T3 & T4 \\
    M1 & 10 & 0 & 3 & 0 \\
    M2 & 0  & 9 & 3 & 0 \\
    M3 & 5  & 5 & 0 & 4 \\
    M4 & 0  & 1 & 3 & 5 \\
  };
  \begin{pgfonlayer}{fondo}
    \marco{a}{5}{5} \marco{b}{5}{5}
    \tfila{a}{2}{5} \tfila{a}{3}{5} \tfila{a}{4}{5} \tfila{a}{5}{5}
  \end{pgfonlayer}
  \cero{b-4-4} \cero{b-5-2} \cero{b-3-5} \cero{b-2-3}
  \flecha{a}{b}{asignar}
  \pie{a}{Paso 3: 4 líneas} \pie{b}{Paso 5: ceros asignados}
\end{tikzpicture}
\end{center}
Con el cero nuevo hacen falta \textbf{4 líneas} (las cuatro filas). Como $4 = 4$, ya se puede asignar.

\paso{Paso 5. Asignar un cero por fila y por columna}
En la tabla de la derecha:
\begin{align*}
1)\quad & \text{M3} \to \text{T3} \nota{único cero de su fila}\\
2)\quad & \text{M4} \to \text{T1} \nota{único cero de su fila; se tacha la columna T1}\\
3)\quad & \text{M2} \to \text{T4} \nota{su cero de T1 ya está tachado: le queda T4}\\
4)\quad & \text{M1} \to \text{T2} \nota{T4 ya está ocupado: le queda T2}
\end{align*}

\subsection*{Resultado}
\[
\res{\text{M1} \to \text{T2},\ \ \text{M2} \to \text{T4},\ \ \text{M3} \to \text{T3},\ \ \text{M4} \to \text{T1},\quad Z^* = 15 \text{ horas}}
\]
\begin{itemize}[itemsep=2pt]
  \item La máquina 1 hace el trabajo 2 (5 h), la 2 el trabajo 4 (5 h), la 3 el trabajo 3 (3 h) y la 4 el trabajo 1 (2 h).
  \item Tiempo total de preparación: $5 + 5 + 3 + 2 = \mathbf{15}$ horas, con los tiempos originales.
  \item La máquina 2 no hace su trabajo más rápido (T1, 2 h): T1 se lo queda la máquina 4, que también tarda 2 h en
        él. Si fuera al revés, M4 tendría que tomar T2 (4 h) y quitárselo a M1, que pasaría a T4 (7 h): el total
        subiría a 16 horas.
\end{itemize}

\textbf{Comprobación:} cada máquina tiene un trabajo y cada trabajo una máquina \checkmark. Lo restado:
filas $5 + 2 + 3 + 2 = 12$, columnas $0 + 0 + 0 + 2 = 2$; el ajuste restó $1$ a dos filas y sumó $1$ a una
columna, o sea bajó $1$. Total $12 + 2 + 1 = 15$ \checkmark. La calculadora revisó además los 24 repartos posibles
y ninguno baja de 15 horas \checkmark.

\subsection*{Interpretación gráfica}
\noindent\begin{minipage}[c]{0.52\textwidth}\centering
\begin{tikzpicture}[>=Stealth]
  \foreach \n/\y in {1/0, 2/-0.9, 3/-1.8, 4/-2.7} \node[ag] (M\n) at (0,\y) {Máquina \n};
  \foreach \n/\y in {1/0, 2/-0.9, 3/-1.8, 4/-2.7} \node[ag] (T\n) at (5,\y) {Trabajo \n};
  \foreach \a in {1,...,4} \foreach \t in {1,...,4}
    \draw[libre] (M\a.east) -- (T\t.west);
  \draw[elegido] (M1.east) -- (T2.west); \node[right=1mm of T2, font=\scriptsize] {5 h};
  \draw[elegido] (M2.east) -- (T4.west); \node[right=1mm of T4, font=\scriptsize] {5 h};
  \draw[elegido] (M3.east) -- (T3.west); \node[right=1mm of T3, font=\scriptsize] {3 h};
  \draw[elegido] (M4.east) -- (T1.west); \node[right=1mm of T1, font=\scriptsize] {2 h};
  \node[font=\small\bfseries, above=2mm of M1] {Máquinas};
  \node[font=\small\bfseries, above=2mm of T1] {Trabajos};
\end{tikzpicture}
\end{minipage}\hfill
\begin{minipage}[c]{0.44\textwidth}
De las 16 parejas posibles se usan 4. En palabras: Machineco debe preparar \textbf{la máquina 1 para el trabajo 2,
la 2 para el trabajo 4, la 3 para el trabajo 3 y la 4 para el trabajo 1}; así pierde solo \textbf{15 horas} en
preparaciones. Este ejercicio muestra por qué a veces hace falta el paso 4: después de restar filas y columnas, dos
máquinas querían el mismo trabajo, y el ajuste con $k = 1$ abrió una segunda opción barata para una de ellas.
\end{minipage}
\par

% =====================================================================
\bigskip
\section*{Resumen de los tres ejercicios}
\begin{center}\small
\begin{tabular}{lccccc}
\toprule
\textbf{Ejercicio} & \textbf{Tamaño} & \textbf{Repartos} & \textbf{Ajustes} & \textbf{Asignación} & \textbf{Óptimo} \\
\midrule
Hijos de Joe Klyne & $3 \times 3$ & 6  & 0 & Pintar, Podar, Lavar & \$27 \\
Fowle Marketing    & $3 \times 3$ & 6  & 1 ($k = 2$) & C2, C3, C1 & 26 días \\
Machineco          & $4 \times 4$ & 24 & 1 ($k = 1$) & T2, T4, T3, T1 & 15 horas \\
\bottomrule
\end{tabular}
\end{center}

\begin{idea}[Conclusiones]
\begin{itemize}[itemsep=2pt, leftmargin=*]
  \item El problema de asignación reparte $n$ tareas entre $n$ agentes, una para cada uno, al menor costo total.
        Probar las $n!$ combinaciones deja de ser posible muy rápido.
  \item Restar una misma cantidad a toda una fila o a toda una columna no cambia cuál es el mejor reparto: solo baja
        el total. El método húngaro usa esto para fabricar ceros.
  \item Si los ceros se tapan con $n$ líneas, se puede elegir un cero por fila y por columna, y ese reparto es óptimo.
        Si no, el ajuste con $k$ crea ceros nuevos.
  \item El costo se calcula siempre con la tabla \textbf{original}, y debe ser igual a la suma de todo lo restado.
        Es un caso especial del transporte y de la programación lineal, pero mucho más cómodo de resolver.
\end{itemize}
\end{idea}

\textbf{Verificación:} los tres ejercicios se comprobaron con la calculadora (\emph{Simplex Arcade}, modo
Asignación), que reproduce los mismos pasos, líneas y valores de $k$; con el modelo de programación lineal resuelto
con su simplex y con la librería javascript-lp-solver; y probando todos los repartos. Todos dan 27, 26 y 15.

\section*{Referencias}
\addcontentsline{toc}{part}{Referencias}
\begin{itemize}[leftmargin=*, itemsep=0pt, topsep=2pt]
  \item Anderson, D. R., Sweeney, D. J. y Williams, T. A. \textit{Métodos cuantitativos para los negocios}. Cengage Learning.
  \item Kuhn, H. W. (1955). The Hungarian method for the assignment problem. \textit{Naval Research Logistics Quarterly}, 2(1--2), 83--97.
  \item Munkres, J. (1957). Algorithms for the assignment and transportation problems. \textit{Journal of the Society for Industrial and Applied Mathematics}, 5(1), 32--38.
  \item Taha, H. A. \textit{Investigación de Operaciones}. Pearson.
  \item Winston, W. L. \textit{Investigación de operaciones: aplicaciones y algoritmos}. Thomson.
\end{itemize}

\end{document}
